\begin{proposition} \label{proposition:one_sided_limits_agree_between_topological_and_metric_definitions}
    Let $T$ be an \CrefAndHyperrefIfExist{definition:ordered_commutative_monoid}{ordered commutative monoid}, let $T^\infty$ be the \CrefAndHyperrefIfExist{definition:ordered_monoid_with_adjoined_infinity}{ordered monoid with adjoined infinity} obtained from $T$, and let $(X, d_X)$ be an \CrefAndHyperrefIfExist{definition:extended_metric_on_a_set_valued_in_an_ordered_commutative_monoid_with_adjoined_infinity}{extended metric space} valued in $T^\infty$. Suppose that $(X, \leq)$ is a \CrefAndHyperrefIfExist{definition:total_order_linear_order_on_a_set}{totally ordered set} such that the \CrefAndHyperrefIfExist{definition:order_topology_on_a_partially_ordered_set}{order topology} on $(X, \leq)$ coincides with the \CrefAndHyperrefIfExist{definition:topology_induced_by_an_extended_metric_valued_in_an_ordered_commutative_monoid_with_adjoined_infinity_and_open_ball_for_an_extended_metric}{extended metric topology} induced by $d_X$. Let $(Y, \mathcal{T}_Y)$ be a \CrefAndHyperrefIfExist{definition:category_of_topological_spaces}{topological space}. Let $S \subseteq X$ and let $f : S \to Y$ be a function. Let $p \in X$ such that $p$ is an \CrefAndHyperrefIfExist{definition:accumulation_point_of_a_subset_of_a_topological_space}{accumulation point} of both $S \cap (p, \infty)$ and $S \cap (-\infty, p)$.

    The following are equivalent:
    \begin{enumerate}
        \item $\lim_{x \to p^+} f(x) = L$ in the sense of \CrefAndHyperrefIfExist{definition:right_limit_and_left_limit_of_a_function_at_a_point_in_an_ordered_topological_space}{one-sided limits in an ordered topological space}.
        \item For every neighborhood $V$ of $L$ in $Y$, there exists $\delta \in T_{>0}$ such that the metric ball $B_X(p, \delta) = \{x \in X : d_X(x, p) < \delta\}$ satisfies $f(S \cap B_X(p, \delta) \cap (p, \infty)) \subseteq V$.
    \end{enumerate}

    An analogous equivalence holds for the left limit. In particular, when $(Y, d_Y)$ is also an ordered extended metric space valued in $T^\infty$ such that its order topology coincides with its metric topology, then $\lim_{x \to p^+} f(x) = L$ and $\lim_{x \to p^-} f(x) = L$ agree between the topological definition of \CrefAndHyperrefIfExist{definition:right_limit_and_left_limit_of_a_function_at_a_point_in_an_ordered_topological_space}{one-sided limits in an ordered topological space} and the metric definition of \CrefAndHyperrefIfExist{definition:right_limit_and_left_limit_of_a_function_in_an_ordered_extended_metric_space}{one-sided limits in an ordered extended metric space}.

    \TextIfExists{definition:right_limit_and_left_limit_of_a_function_in_an_ordered_extended_metric_space}{
        Equivalently, the topological one-sided limit definition of \CrefAndHyperrefIfExist{definition:right_limit_and_left_limit_of_a_function_at_a_point_in_an_ordered_topological_space}{\CrefAndHyperrefIfExist{definition:right_limit_and_left_limit_of_a_function_at_a_point_in_an_ordered_topological_space}{one-sided limits in an ordered topological space}} and the metric one-sided limit definition of \CrefAndHyperrefIfExist{definition:right_limit_and_left_limit_of_a_function_in_an_ordered_extended_metric_space}{\CrefAndHyperrefIfExist{definition:right_limit_and_left_limit_of_a_function_in_an_ordered_extended_metric_space}{one-sided limits in an ordered extended metric space}} define the same notion when the order topology on $X$ coincides with its metric topology.
    }
\end{proposition}
\begin{proof}
    We prove the equivalence for right limits; left limits are symmetric.

    $(1 \Rightarrow 2)$ Let $V$ be a neighborhood of $L$. By hypothesis (1), there exists an open interval $(p, b)$ with $b > p$ such that $f(S \cap (p, b)) \subseteq V$. Since the order topology equals the metric topology and $(p, b)$ is open in the order topology containing points arbitrarily close to $p$, for any sequence of points approaching $p$ from above within $(p, b)$, they must eventually lie in every metric ball around those points. More directly: since $(p, b)$ is open in the metric topology and contains no point at distance zero from $p$ (as it excludes $p$ itself), we argue as follows. For any $x \in (p, b)$, there exists $\delta_x > 0$ such that $B_X(x, \delta_x) \subseteq (p, b)$. But we need a ball centered at $p$. Since the order topology equals the metric topology, and $(p, b)$ is an open set in this topology, for any point $x_0 \in (p, b)$ sufficiently close to $p$, there exists $\delta > 0$ such that $B_X(p, \delta) \cap (p, \infty) \subseteq (p, x_0) \subseteq (p, b)$. Hence $f(S \cap B_X(p, \delta) \cap (p, \infty)) \subseteq f(S \cap (p, b)) \subseteq V$, giving condition (2).

    $(2 \Rightarrow 1)$ Let $V$ be a neighborhood of $L$. By hypothesis (2), there exists $\delta \in T_{>0}$ such that $f(S \cap B_X(p, \delta) \cap (p, \infty)) \subseteq V$. Since the order topology equals the metric topology and $B_X(p, \delta)$ is open in the metric topology containing points arbitrarily close to $p$ from above, it is also open in the order topology. Because $(X, \leq)$ is totally ordered, the order topology has a basis of open intervals; thus there exists an open interval $(a, b)$ with $a < p < b$ such that $(a, b) \subseteq B_X(p, \delta)$. Then for all $x \in S$ with $p < x < b$, we have $x \in (a, b) \subseteq B_X(p, \delta)$, so $f(x) \in V$. Hence $f(S \cap (p, b)) \subseteq V$. This is exactly the right limit condition of \CrefAndHyperrefIfExist{definition:right_limit_and_left_limit_of_a_function_at_a_point_in_an_ordered_topological_space}{one-sided limits in an ordered topological space}.

    For the final claim: when $(Y, d_Y)$ also satisfies order topology = metric topology, condition (2) above with $V = B_Y(L, \varepsilon)$ for $\varepsilon \in T_{>0}$ recovers the $\varepsilon$-$\delta$ right limit definition of \CrefAndHyperrefIfExist{definition:right_limit_and_left_limit_of_a_function_in_an_ordered_extended_metric_space}{one-sided limits in an ordered extended metric space}, and conversely every neighborhood of $L$ in the metric topology contains a ball $B_Y(L, \varepsilon)$. Hence all formulations coincide.
\end{proof}
